\chapter{Example chapter}
% Give your sections labels so you can refer to them from other spots (this works across .tex files in the "sections" folder too!).
\label{ch:example}

% Include a statement of originality describing your contributions at the start of the chapter
\begin{originality}
	For the work presented in this chapter I have contributed <X>.
	Most figures in this section, as well as parts of the relevant discussion, are adapted from the following publication, plotting for which has been done by me, and a couple of other people:

	J.~Dranczewski,\footnotemark[1] W.~K.~Ng,\footnotemark[1] A.~Fischer,\footnotemark[1] T.~V.~Raziman, D.~Saxena, T.~Farchy, K.~Stenning, J.~Peters, H.~Schmid, W.~R.~Branford, M.~Barahona, K.~Moselund, R.~Sapienza, J.~C.~Gartside, ``\href{https://doi.org/10.48550/arXiv.2407.15558}{Retinomorphic Machine Vision in a Network Laser},'' \textit{arXiv}, Aug. 2024. \textit{[manuscript in review]}
\end{originality}

%\footnotetext[1]{Authors contributed equally.}
% Use a custom text instead of the default footnote, which would go to the bottom of the page.
\vspace{-1.5em}
{\footnotesize*Authors contributed equally.}

\section*{Abstract}
% Include a brief chapter abstract if desired.
The header images are optional, and described in more detail in `thesis\_style.sty'. Basically, put a .pdf file matching the chapter number/letter in `figures/\_headers`, and you'll get a chapter header image. If there's no file for a chapter, it defaults to `default.pdf`. Remove that file -- no headers!

\lipsum[1]

\section{Example section}
\label{sec:example_section}
\begin{outline}
	\item You can write notes to yourself in ``outline'' environments. They act like LaTeX's ``itemize'', but are coloured blue, and you can easily hide them from the final document once you're done with a specific outline by changing the environment name to ``outlineDone''.
	\item I also provide a number of potentially useful icons that can be used to annotate these outlines:
	\begin{itemize}
		\item Checkmarks, for points you've covered \checkmark
		\item Hourglass for things that are in progress \progress
		\item Paper, for things to be adapted from papers \paper
		\item Globe, for things that need to be referenced or researched from external sources \globe
		\item Notion, as that's where I was tracking figure progress \notion
	\end{itemize}
	\item You can use these anywhere in the text just by using the commands listed in the source above: \checkmark\paper
\end{outline}

\begin{outlineDone}
	\item This outline will not appear in the compiled document, but you can keep it in the source for reference.
\end{outlineDone}

Example citation\cite{dranczewskiPlasmaEtchingFabrication2023} -- citations appear in the upper right by default, if you need to reference a citation in the main text you can use Ref.~\citein{dranczewskiPlasmaEtchingFabrication2023}. Note the \goodtilde~symbol -- it forms a non-breaking space between ``Ref.'' and the citation number, meaning that they will always appear on the same line. You can similarly reference chapters and sections, for example Chapter~\ref{ch:example}.

Here's an example figure: Fig.~\ref{fig:z_example_a_name}

\begin{figure}[h]
	\centering
	% You'll note that I'm not including a figure width here - in my opinion you should prepare the figures
	% to be the correct width from the start, which avoids any rescaling issues or font size discrepancies.
	% For this document, the figure width should be 6.3 inches
	\includegraphics{z_example_a_name/z_example_a_name.pdf}
	\caption[
		Wide example figure % This short description will appear in the figure list
	]{
		% This long description will appear in the text.
		Example figure description.
		\textbf{a}~This describes part \textit{a} of the figure.
		\textbf{b}~This describes part \textit{b} of the figure.
	}
	% Provide an ID for the figure so it can be referenced elsewhere in the text.
	% Mine follow a scheme of [sequential chapter letter]_[short chapter name]_[sequential figure letter]_[short figure name]
	% This allows easy alphabetical sorting of the figure folders, and general consistency.
	\label{fig:z_example_a_name}
\end{figure}

Here's an example of a narrow figure with the caption to the side: Fig.~\ref{fig:a_example_b_narrow}. You can use `figures/placeholder.pdf' as temporary content while you work on the real figures.

% Here we make the figure narrower, add some horizontal fill, and then add the caption to the right in a "minipage"
\begin{figure}[h]
	% If you do need to specify a figure width, this is how.
	
\includegraphics[width=2.2in, valign=c]{placeholder.pdf}
	\hfill
	\begin{minipage}{\dimexpr \textwidth-\columnsep-2.2in}
		\captionsetup{skip=0pt}
		\caption[
			Narrow example figure
		]{
			Narrow figure caption.
		}
		\label{fig:a_example_b_narrow}
	\end{minipage}
\end{figure}

Examples of acronym usage: first instance -- \ac{CPU}, second instance -- \ac{CPU}. Forced full: \acf{CPU}.
Examples of `siunitx' usage: \qty{1.23}{\meter\squared}, \qty{45.6}{\fluence} (custom unit defined in `thesis\_style.tex').
Chemical formula: \ch{Al2(SO4)3}